\section{Embedding into the JTS Forcing Translation: the Propositional Fragment}
\label{sec:bridge}

A parallel literature, alongside the topos-of-trees account of
step-indexing surveyed in Section~\ref{sec:topos}, develops
\emph{syntactic} forcing translations of type theories.  Jaber,
Tabareau, and Sozeau \cite{jaber-tabareau-sozeau-2012} give a forcing
translation of the Calculus of Constructions: for any forcing notion
$P$, they produce a syntactic translation that takes a CIC term
$M : T$ to a forced term $M^P : T^P$, where the type translation
$T \mapsto T^P$ corresponds to ``the $P$-forced version of $T$''.  At
the level of propositions, the forced version of $\mathsf{Prop}$ is
essentially a downward-closed presheaf on the forcing conditions
\cite{jaber-et-al-2016}, which is the type of our $\mathtt{iProp}$.

The framework of Section~\ref{sec:abstract} is therefore, up to
definitional choices, the propositional fragment of the JTS forcing
translation of CIC, instantiated at an arbitrary forcing structure.
At the propositional level, this identification is now a
machine-checked theorem rather than an observation
(\texttt{Phase6b\_JTS.v}).  We deep-embed a propositional object
language \texttt{form} (atoms, $\top$, $\bot$, $\wedge$, $\vee$,
$\to$, $\triangleright$) and define the forcing translation
\texttt{jts} exactly as the JTS clauses prescribe --- a standalone
\texttt{Fixpoint} into \texttt{cond -> Prop} that never mentions
\texttt{iProp}: implication quantifies over all future conditions,
$\triangleright$ over all strict predecessors, the other connectives
are read pointwise.  Three theorems, all over an \emph{arbitrary}
forcing structure, cash out the claim:

\begin{lstlisting}
Lemma   jts_mono      : (* translated formulas are presheaves:  *)
  forall f p q, q ≼ p -> jts f p -> jts f q.
Theorem jts_is_interp : (* translation = iProp semantics,       *)
  forall f p,           (* pointwise, at every formula          *)
    jts f p <-> ihold FS (interp f) p.
Theorem jts_Lob       : (* Löb holds OF the translated formulas *)
  forall f p, jts (FImpl (FLater f) f) p -> jts f p.
\end{lstlisting}

\noindent
Here \texttt{interp} reads the same syntax through the semantic
connectives of Section~\ref{sec:abstract}.  \texttt{jts\_mono} is the
presheaf condition (the translation lands in \texttt{iProp});
\texttt{jts\_is\_interp} is the factorization ``syntactic forcing
translation $=$ step-indexed semantics''; and \texttt{jts\_Lob} is
the step-indexed dividend, transported to the syntax --- available
precisely because the forcing structure is well-founded
(Section~\ref{sec:lob-essential} shows it is otherwise unavailable).
What remains unmechanized is the \emph{term-level} translation of
full CIC (universes, dependent products, the definitional side of
\cite{jaber-et-al-2016}); the propositional claim made above is
mechanized in full.

Alongside this, we mechanize the \emph{embedding} of the unforced
theory ($\mathsf{Prop}$ in the ambient meta-theory) into the forced
one: the constant embedding.  It makes precise what the forcing
translation does to a meta-level proposition that is to be lifted
into the forced logic unchanged.

\paragraph{Relation to the available JTS Coq plugin.}
The JTS translation has been implemented as a Coq plugin
\cite{forcing-coq-plugin}.  It is worth being explicit about how it
differs from the work of this paper, because the existence of a
``forcing in Coq'' plugin invites the question of overlap.  The
plugin is a \emph{translation tool}: it takes a Coq term $M : T$
together with a forcing structure, and produces a translated term
$M^P : T^P$ in the forced universe.  Its main demonstrative example
is a step-indexed model of pure $\lambda$-calculus, in which the
translation lets one define a $\mathtt{fixp}$ operator and prove
equations such as $Y$-combinator reduction. The relevant
infrastructure lives in the test-suite files \texttt{fix.v} and
\texttt{pure\_lambda\_calculus.v}.  The instantiation at the natural
numbers with $\le$ (\texttt{NatForcing} in the plugin's \texttt{Init.v})
produces a structure of monotone-indexed sheaves which is, up to
definitional details, exactly the type of our $\mathtt{iProp}$.

The plugin does not contain a Hoare logic, a weakest-precondition
predicate, an operational semantics for any source language, or any
soundness proof against an operational semantics.  In particular it
does not state any rule of the shape $(\triangleright \varphi \to
\varphi) \vdash \varphi$ as the soundness justification for a
recursive program-logic rule.  The work of this paper is therefore
disjoint from the plugin's contribution: we use the same forcing
structure (and the same propositional fragment of the translation)
that the plugin makes available, but we apply it to a different
purpose: a step-indexed Hoare logic, the diagnostic of
Section~\ref{sec:imp-vs-lam}, and the well-foundedness counterexample
of Section~\ref{sec:lob-essential}.  The plugin would be a natural
substrate to build on if one wished to formalize the
\emph{type-level} half of the correspondence we sketch here as
future work.

\paragraph{The constant embedding.}
The embedding $\mathsf{Prop} \hookrightarrow \mathtt{iProp}$ sends a
meta-level proposition $P$ to the iProp that is forced everywhere if
$P$ is true and forced nowhere if $P$ is false.  Either way it is
trivially downward-closed in the forcing order, since it does not
depend on the condition:

\begin{lstlisting}
Definition embed (P : Prop) : iProp :=
  MkProp (fun _ => P) (embed_mono P).

Notation "⌜ P ⌝" := (embed P).
\end{lstlisting}

\noindent
This is the trivial part of the forcing translation: classical
propositional reasoning lifts to the forcing framework without
change.  We will see in a moment that the genuinely new content of
the forcing translation, the $\triangleright$ modality, is
\emph{not} in the image of $\lceil \cdot \rceil$.

\paragraph{Commutation with the Heyting connectives.}
The embedding is a Heyting algebra homomorphism: it commutes with
each of $\top$, $\bot$, $\wedge$, $\vee$, $\to$ in both directions of
entailment.  The proofs are short and structural:

\begin{lstlisting}
Lemma embed_True_l : ⌜ True ⌝ ⊢ iTrue.
Lemma embed_True_r : iTrue ⊢ ⌜ True ⌝.
Lemma embed_False_l : ⌜ False ⌝ ⊢ iFalse.
Lemma embed_False_r : iFalse ⊢ ⌜ False ⌝.
Lemma embed_and_intro P Q : ⌜ P /\ Q ⌝ ⊢ ⌜ P ⌝ ∧ ⌜ Q ⌝.
Lemma embed_and_elim P Q : ⌜ P ⌝ ∧ ⌜ Q ⌝ ⊢ ⌜ P /\ Q ⌝.
Lemma embed_or_intro P Q : ⌜ P \/ Q ⌝ ⊢ ⌜ P ⌝ ∨ ⌜ Q ⌝.
Lemma embed_or_elim P Q : ⌜ P ⌝ ∨ ⌜ Q ⌝ ⊢ ⌜ P \/ Q ⌝.
Lemma embed_impl_intro P Q : ⌜ P -> Q ⌝ ⊢ (⌜ P ⌝ → ⌜ Q ⌝).
Lemma embed_impl_elim P Q : (⌜ P ⌝ → ⌜ Q ⌝) ⊢ ⌜ P -> Q ⌝.
\end{lstlisting}

\noindent
The most interesting case is the implication direction
$(\lceil P \rceil \to \lceil Q \rceil) \vdash \lceil P \to Q \rceil$,
where the Kripke implication of $\to$ has to be discharged by an
appeal to the reflexive case ``$n \le n$''.  Given an inhabitant of
the Kripke implication at depth $n$, we apply it at $n$ itself, with
the witness of $n \le n$ given by \texttt{Nat.le\_refl}, and obtain
the meta-level implication.  The other cases are direct unfoldings
of the connectives' Kripke definitions.

\paragraph{Why this matters.}
Each commutation lemma corresponds to a clause of the propositional
fragment of the JTS forcing translation.  Specifically, the
translation $\llbracket P \rrbracket^P$ at the level of propositions
commutes with the standard connectives in exactly the way our
$\lceil \cdot \rceil$ does: a translated conjunction is the
conjunction of translations, and so on \cite{jaber-tabareau-sozeau-2012}.
Our framework therefore captures the propositional fragment of the
translation.  What we do not formalize is the dependent fragment, in
which the translation acts non-trivially on universally quantified
types and on the universes themselves.  That is the part of the
correspondence which would, if mechanized, turn the present
informal-correspondence claim into a definitional-equality theorem.

\paragraph{The later modality is new content.}
The constant embedding is not surjective.  In particular, its image
is not closed under $\triangleright$: there is a meta-level
proposition $P$ such that $\triangleright \lceil P \rceil$ is not
entailment-equivalent to $\lceil P \rceil$.  We exhibit the
witness:

\begin{lstlisting}
Lemma embed_later_distinct :
  ~ (▷ ⌜ False ⌝ ⊢ ⌜ False ⌝).
Proof.
  intros H. apply (H 0). simpl. trivial.
Qed.
\end{lstlisting}

\noindent
At depth $0$, $\triangleright \lceil \mathrm{False} \rceil$ is
vacuously \texttt{True}, whereas $\lceil \mathrm{False} \rceil$ at
depth $0$ is the constant \texttt{False}.  The entailment from one to
the other therefore fails at depth $0$, which is exactly the
witness that $\triangleright$ produces a proposition outside the
embedded image.

This makes precise what the step-indexed-specific content of our
framework is.  The Heyting algebra structure of $\mathtt{iProp}$,
all the way through $\top, \bot, \wedge, \vee, \to$, is
captured by the embedded image of $\mathsf{Prop}$.  What the
step-indexed framework adds is the $\triangleright$ modality, and
the rules that come with it (monotonicity, $\varphi \vdash
\triangleright \varphi$, distribution over $\wedge$, and Löb's
rule).  In the JTS picture, this is the structure that comes from
the forcing notion specifically, rather than from the structural
machinery of the type theory. In our picture, it is the structure
that comes from the well-founded strict refinement of the forcing
conditions, rather than from the constant embedding.

\paragraph{A summary correspondence.}
The picture that emerges is this.  Step-indexed Hoare logic, as
practiced in Iris and elsewhere, has two layers of structure.  The
first is a Heyting-algebra fragment, which lifts from the meta-theory
by a constant embedding and therefore has no forcing-specific content.
The second is a modal fragment ($\triangleright$, $\mathtt{iLob}$, and
the rules around them), which is specifically the structure that the
well-founded forcing notion contributes.  The JTS forcing
translation is the type-theoretic analog of this picture: for any
forcing notion $P$, it gives a translation of CIC that adds
$P$-specific structure to the embedded image of the unforced
theory.  At the level of propositions, the correspondence we
formalize is the constant embedding.  A full correspondence at all
levels of the universe would require mechanizing the JTS
translation itself, which we leave as future work.

